\sect{पक्षवर्धिनी-एकादशी-माहात्म्यम्}
\label{sec:padma-pakshavardhini}

\ganapatyadiDhyanam
\padmaPuranam

\dnsub{अथ नामाष्टत्रिंशोऽध्यायः॥३८}

\uvacha{नारद उवाच}

\twolineshloka
{कीदृशी स्यान्महादेव पक्षवर्धिनिसंज्ञका}
{यया वै कृतया जन्तुर्महपापात् प्रमुच्यते}% ॥१॥

\uvacha{श्रीमहादेव उवाच}

\onelineshloka*
{अमा वा यदि वा पूर्णा सम्पूर्णा जायते तदा}

\twolineshloka
{भूत्वा वै नाडिकाषष्ठिर्जायते प्रतिपद्दिने}
{अश्वमेधायुतैस्तुल्या सा भवेत् पक्षवर्धिनी}% ॥२॥

\uvacha{नारद उवाच}

\twolineshloka
{पूजाविधिं तु पृच्छामि साम्प्रतं देवसत्तम}
{यत्कृते तु महादेव महाफलमवाप्नुयात्}% ॥३॥

\uvacha{महादेव उवाच}

\twolineshloka
{पूजाविधिं प्रवक्ष्यामि साम्प्रतं द्विजनन्दन}
{पूजिते चार्चिते विष्णौ फलं प्राप्नोत्यसंशयः}% ॥४॥

\twolineshloka
{येन पूजाविधानेन तुष्टिं प्राप्नोति माधवः}
{अव्रणं जलपूर्णं च कुम्भं चन्दनचर्चितम्}% ॥५॥

\twolineshloka
{पञ्चरत्नसमायुक्तं पुष्पमालाभिवेष्टितम्}
{स्थाप्यं ताम्रमयं पात्रं सगोधूमं घटोपरि}% ॥६॥

\twolineshloka
{सौवर्णं कारयेद् देवं माससंज्ञाभिनामकम्}
{पञ्चामृतेन विधिना स्नपनं सुमनोरमम्}% ॥७॥

\twolineshloka
{कारयेद् देवदेवेशं जगन्नाथं जगत्पतिम्}
{विलेपनं तु कर्तव्यं कुङ्कुमागरुचन्दनैः}% ॥८॥

\twolineshloka
{वस्त्रयुग्मं च दातव्यं छत्रोपानहसंयुतम्}
{पूजयेद् देवताधीशं कुम्भपात्रोपरिस्थितम्}% ॥९॥

\twolineshloka
{पद्मनाभाय वै पादौ जानुनी विश्वमूर्तये}
{ऊरू वै ज्ञानगम्याय कटी दानप्रदाय च}% ॥१०॥

\twolineshloka
{उदरं विश्वनाथाय हृदयं श्रीधराय च}
{कण्ठं कौस्तुभकण्ठाय बाहू क्षत्रान्तकारिणे}% ॥११॥

\twolineshloka
{ललाटं व्योममूर्ध्ने तु शिरो वै सर्वरूपिणे}
{स्वनाम्ना चैव कमलां सर्वाङ्गीं दिव्यरूपिणीम्}% ॥१२॥

\twolineshloka
{एवं विधिवत्सम्पूज्य ततोऽर्घं दापयेत्सुधीः}
{नालिकेरेण शुभ्रेण देवदेवस्य चक्रिणः}% ॥१३॥

\twolineshloka
{अनेनैवार्घदानेन सम्पूर्णं जायते व्रतम्}
{संसारार्णवमग्नं मां समुद्धर जगत्पते}% ॥१४॥

\twolineshloka
{त्वमीशः सर्वलोकानां त्वं साक्षाच्च जगत्पतिः}
{गृहाणार्घं मया दत्तं पद्मनाभ नमोस्तु ते}% ॥१५॥

\twolineshloka
{नैवेद्यानि सुहृद्यानि षड्रसानि विशेषतः}
{देयानि तु विशेषेण केशवाय सुभक्तितः}% ॥१६॥

\twolineshloka
{नागपत्रं सकर्पूरं दद्याद् देवस्य भक्तितः}
{घृतेन दीपकं कुर्यात्तिलतैलेन वा पुनः}% ॥१७॥

\twolineshloka
{कृत्वा सम्यग्विधानेन गुरोः पूजां प्रकारयेत्}
{वस्त्राणि चैव चोष्णीषं कञ्चुकं च प्रदापयेत्}% ॥१८॥

\twolineshloka
{दक्षिणां च यथाशक्त्या गुरुवे सम्प्रदापयेत् }
{भोजनं चैव ताम्बूलं दत्त्वा चार्घं प्रदापयेत्}% ॥१९॥

\twolineshloka
{स्ववित्तस्यानुमानेन यथाशक्त्या तु निर्धनैः}
{कार्या सम्यक्प्रयत्नेन द्वादशी पक्षवर्धिनी}% ॥२०॥

\twolineshloka
{ततो जागरणं कुर्याद्गीतनृत्यसमन्वितम्}
{पुराणपाठसहितं हास्याह्लादसमन्वितम्}% ॥२१॥

\twolineshloka
{स्तुवन्ति च प्रशंसन्ति जागरं चक्रपाणिनः}
{नित्योत्सवो भवेत्तेषां गृहे वै दशजन्मसु}% ॥२२॥

\twolineshloka
{अतो धन्यतमा चेयं कर्तव्या पक्षवर्धनी}
{कृत्वा तु सकलं पुण्यं फलं प्राप्नोत्यसंशयम्}% ॥२३॥

\twolineshloka
{पक्षवर्धिनी-माहात्म्यं ये शृण्वन्ति मनीषिणः}
{तैः कृतं सत्कृतं सर्वं यावदाभूतसम्प्लवम्}% ॥२४॥

\twolineshloka
{पञ्चाग्निसाधने पुण्यं यच्च स्यात्तीर्थसाधने}
{तत्पुण्यं समवाप्नोति विष्णोर्जागरकारणात्}% ॥२५॥

\twolineshloka
{पक्षवर्धनिका पुण्या पवित्रा पापनाशिनी}
{उपवासकृता विप्र हत्याकोटिविनाशिनी}% ॥२६॥

\twolineshloka
{वसिष्ठेन कृता पूर्वं भारद्वाजेन वै मुने}
{ध्रुवेण चाम्बरीषेण कृतेयं विष्णुवल्लभा}% ॥२७॥

\twolineshloka
{इयं काशीसमा पुण्या इयं वै द्वारका समा}
{उपोषिता च भक्तेन वाञ्छितं च ददात्यसौ}% ॥२८॥

\twolineshloka
{इयं धन्या धन्यतमा हत्यायुतविनाशिनी}
{कर्तव्या तु विशेषेण वैष्णवैर्ज्ञानतत्परैः}% ॥२९॥

\twolineshloka
{अहो सर्वेश्वरो देवः संसेव्यो व्रततत्परैः}
{किमन्यद् बहुनोक्तेन कर्तव्यं व्रतमुत्तमम्}% ॥३०॥

\twolineshloka
{यथा च वर्धते चन्द्रः सिते पक्षे विशेषतः}
{तथा वै वर्धते भक्त कारणात् पक्षवर्धिनी}% ॥३१॥

\twolineshloka
{सूर्योदये यथा ध्वान्तं नश्यते तत्क्षणादपि}
{तथाऽघं नाशमाप्नोति करणात् पक्षवर्धिनी}% ॥३२॥

॥इति श्रीपाद्मे महापुराणे पञ्चपञ्चाशत्साहस्र्यां संहितायामुत्तरखण्डे उमापतिनारदसंवादान्तर्गतकृष्णयुधिष्ठिरसंवादे पक्षवर्धिनी-एकादशी-माहात्म्यं नाम नामाष्टत्रिंशोऽध्यायः॥३८॥

\genMangalaShloka

\hyperref[sec:ekadashi_mahatmyam_padma_puranam]{\closesub}
\clearpage

